% Fragment for the consolidated manuscript: TAUM's moving decompositions, projector
% curvature and characteristic classes (TAUM lines 2008-2508).
% Body only. Build with fragments/build_fragment.py moving_subspace.
\section{When the selected subspace moves}
\label{sec:moving-subspace}

The sections so far fix one map $A$ and the splitting into an input and an output sector. Many realizations also supply a family of subspaces. In the pair model the prepared plane lies inside the output space and depends on the setting $\theta$ of the preparation and on the time waited afterwards, while the splitting into the two sectors stays fixed. This section treats a subspace of a fixed Hilbert space that depends differentiably on parameters. Three questions arise. What is the derivative of a moving subspace? How should a vector be carried from one subspace of the family to a nearby one, and does the result depend on the route? Which quantities attached to the family are integers?

The answers use only the family of orthogonal projections. The derivative in one direction is a bounded map from the subspace to its complement, together with its adjoint, so each point and direction supplies an instance of the structure of Section~\ref{sec:cohomology-index}. Carrying a vector along a path is parallel transport. Its local dependence on the route is measured by a curvature, which is the antisymmetric part of the product of two such maps. Around a loop, transport returns the subspace to itself by a unitary, the holonomy, which can be nontrivial even where the family restricted to the loop has no curvature. For subspaces of finite dimension the trace of the curvature, integrated over a closed surface of parameters and multiplied by $i/(2\pi)$, is an integer.

The parameter space and the family of projections are supplied data. Nothing in this section makes the resulting connection a dynamical field. It has no field equation, and the parameters need not be points of spacetime.

\subsection{The derivative of a projection}
\label{subsec:moving-derivative}

Let $\mathcal K$ be a complex Hilbert space, $X$ a finite-dimensional manifold of class $C^2$, and
\begin{equation}
 \Pi:X\to\mathcal B(\mathcal K),\qquad \Pi(x)^2=\Pi(x)=\Pi(x)^* ,
 \label{eq:moving-projection-field}
\end{equation}
a map that is twice continuously differentiable in operator norm. Write $\Pi^\perp=I-\Pi$, $\mathcal L_x=\Ran\Pi(x)$ and $\mathcal L_x^\perp=\Ran\Pi^\perp(x)$. Each point carries the symmetry
\begin{equation}
 J_\Pi(x)=2\Pi(x)-I=\Pi(x)-\Pi^\perp(x),
 \label{eq:moving-symmetry}
\end{equation}
which plays the part of $J$ in Section~\ref{sec:neutral-graphs-phase} for the splitting $\mathcal K=\mathcal L_x\oplus\mathcal L_x^\perp$. The Hilbert inner product of $\mathcal K$ is the same at every point. The form $[z,w]_x=\langle z,J_\Pi(x)w\rangle$ varies with $x$. The space $\mathcal K$ may be the direct sum of the two sectors, or a space inside one of them. The fixed splitting of the earlier sections is the constant field equal to the projection onto the input sector, for which $J_\Pi$ is the symmetry $J$ of Section~\ref{sec:cohomology-index}.

\begin{proposition}[Nearby projections are unitarily equivalent]
\label{prop:moving-unitary}
Let $P,Q$ be orthogonal projections on $\mathcal K$ with $\|P-Q\|<1$, and put $R=(P-Q)^2$. Then
\begin{equation}
 T=\bigl[QP+(I-Q)(I-P)\bigr](I-R)^{-1/2}
 \label{eq:moving-direct-rotation}
\end{equation}
is unitary and $TPT^*=Q$. Consequently $\dim\mathcal L_x$ and $\dim\mathcal L_x^\perp$ are constant on each connected component of $X$. Near each $x_0\in X$ the unitaries $T(x)$ formed from $P=\Pi(x_0)$ and $Q=\Pi(x)$ are of class $C^2$ in $x$ and carry $\mathcal L_{x_0}$ onto $\mathcal L_x$.
\end{proposition}
\begin{proof}
Expanding gives $R=P+Q-PQ-QP$ and $PR=P-PQP=RP$, and likewise $QR=RQ$. Put $T_0=QP+(I-Q)(I-P)$. Then
\[
 T_0^*T_0=PQP+(I-P)(I-Q)(I-P)=I-R ,
\]
and in the same way $T_0T_0^*=I-R$. Since $\|R\|<1$, the operator $I-R$ is positive and invertible, and it commutes with $T_0$. Hence $T=T_0(I-R)^{-1/2}$ satisfies $T^*T=TT^*=I$. Finally $T_0P=QP=QT_0$ gives $TP=QT$. This is the direct rotation of a pair of projections \cite[Ch.~I, \S\S4.6 and 6.8]{Kato}. Unitary equivalence preserves both dimensions, and the set where $\|\Pi(x)-\Pi(x_0)\|<1$ is open, so the dimensions are locally constant. The inverse square root is a norm-convergent power series in $R$, so $T(x)$ is as differentiable as $\Pi(x)$.
\end{proof}

The subspaces $\mathcal L_x$ therefore form a bundle over each component of $X$, of finite or infinite rank, which the unitaries $T(x)$ trivialize locally.

\begin{proposition}[Tangent map of a moving subspace]
\label{prop:moving-tangent}
Let $Y$ be a vector field of class $C^1$ on $X$, and write $Y\Pi$ for the derivative of $\Pi$ along $Y$. Then $Y\Pi$ is self-adjoint,
\begin{equation}
 \Pi\,(Y\Pi)\,\Pi=0,\qquad \Pi^\perp(Y\Pi)\,\Pi^\perp=0,
 \label{eq:moving-diagonal-zero}
\end{equation}
and $J_\Pi\,(Y\Pi)=-(Y\Pi)\,J_\Pi$. Define the tangent map
\begin{equation}
 B_Y(x)=\Pi^\perp(x)\,(Y\Pi)(x)\,\Pi(x):\mathcal L_x\to\mathcal L_x^\perp,\qquad
 B_Y(x)^*=\Pi(x)\,(Y\Pi)(x)\,\Pi^\perp(x).
 \label{eq:moving-tangent-map}
\end{equation}
Relative to $\mathcal K=\mathcal L_x\oplus\mathcal L_x^\perp$,
\begin{equation}
 Y\Pi=\begin{pmatrix}0&B_Y^*\\B_Y&0\end{pmatrix},\qquad
 J_\Pi\,Y\Pi=\begin{pmatrix}0&B_Y^*\\-B_Y&0\end{pmatrix},
 \label{eq:moving-block-forms}
\end{equation}
and
\begin{equation}
 (Y\Pi)^2=\begin{pmatrix}B_Y^*B_Y&0\\0&B_YB_Y^*\end{pmatrix}=-(J_\Pi\,Y\Pi)^2 .
 \label{eq:moving-squares}
\end{equation}
\end{proposition}
\begin{proof}
Differentiating $\Pi^2=\Pi$ along $Y$ gives $(Y\Pi)\Pi+\Pi(Y\Pi)=Y\Pi$. Multiplying by $\Pi$ on both sides gives $2\Pi(Y\Pi)\Pi=\Pi(Y\Pi)\Pi$, the first identity. The field $\Pi^\perp$ consists of projections and $Y\Pi^\perp=-Y\Pi$, which gives the second. The derivative of a self-adjoint field is self-adjoint. With $J_\Pi=2\Pi-I$,
\[
 J_\Pi(Y\Pi)+(Y\Pi)J_\Pi=2\bigl[\Pi(Y\Pi)+(Y\Pi)\Pi\bigr]-2\,Y\Pi=0 .
\]
Inserting $I=\Pi+\Pi^\perp$ on both sides of $Y\Pi$ and using Eq.~\eqref{eq:moving-diagonal-zero} leaves the two off-diagonal blocks. Multiplication by $J_\Pi=\operatorname{diag}(I,-I)$ changes the sign of the lower row. The squares follow by multiplying the block forms, with $J_\Pi^2=I$ and the anticommutation.
\end{proof}

The map $B_Y(x)$ is bounded, depends only on the value $Y(x)$, and is real-linear in it. In coordinates $x^\mu$ write $B_\mu=B_{\partial_\mu}=\Pi^\perp(\partial_\mu\Pi)\Pi$.

At a point $x$ and for a direction $Y$, the spaces $\mathcal L_x$ and $\mathcal L_x^\perp$ with the map $B_Y(x)$ are an instance of the two-term complex of Section~\ref{sec:cohomology-index}, with $\mathcal H_\Psi$, $\mathcal H_\Phi$ and $A$ replaced by $\mathcal L_x$, $\mathcal L_x^\perp$ and $B_Y(x)$. Its differential, regarded as an operator on $\mathcal K$, is $d_Y=\Pi^\perp(Y\Pi)\Pi$, and
\begin{equation}
 Y\Pi=d_Y+d_Y^*,\qquad J_\Pi\,Y\Pi=d_Y^*-d_Y,\qquad (Y\Pi)^2=d_Y^*d_Y+d_Yd_Y^* .
 \label{eq:moving-two-odd}
\end{equation}
These are the two odd operators named after Eq.~\eqref{eq:index-graded} and the Laplacian of that section. The derivative $Y\Pi$ is Hilbert-self-adjoint. It is the tangent vector to the family of subspaces. Multiplying it by the symmetry gives $J_\Pi\,Y\Pi$, which is Hilbert-skew-adjoint and is self-adjoint for the form $[\cdot,\cdot]_x$, in the sense $T^\sharp=J_\Pi T^*J_\Pi$ of Section~\ref{sec:neutral-graphs-phase}. The results of Sections~\ref{sec:polar-matched-observables}, \ref{sec:neutral-graphs-phase} and~\ref{sec:cohomology-index} that concern a single bounded map therefore apply to each $B_Y(x)$: its polar decomposition, the pairing of the nonzero spectra of $B_Y^*B_Y$ and $B_YB_Y^*$, and its index when it is Fredholm. A moving subspace adds no new algebra at one point and in one direction. What it adds is a relation between different directions.

\subsection{The projected connection and its curvature}
\label{subsec:moving-connection}

A section is a map $s:X\to\mathcal K$ with $s(x)\in\mathcal L_x$ for every $x$, that is, $\Pi s=s$. The derivative $Ys$ of a section, taken in the fixed space $\mathcal K$, need not lie in $\mathcal L_x$. The projected connection, also called the Grassmann connection, keeps the part that does:
\begin{equation}
 \nabla_Ys=\Pi\,(Ys).
 \label{eq:moving-connection}
\end{equation}
It is compatible with the inner product. For sections $s,s'$ of class $C^1$, $Y\langle s,s'\rangle=\langle\nabla_Ys,s'\rangle+\langle s,\nabla_Ys'\rangle$, because $\Pi s=s$ and $\Pi s'=s'$.

\begin{proposition}[Curvature of the projected connection]
\label{prop:moving-curvature}
Let $Y,Z$ be vector fields of class $C^1$ and $s$ a section of class $C^2$. Then
\begin{equation}
 \bigl(\nabla_Y\nabla_Z-\nabla_Z\nabla_Y-\nabla_{[Y,Z]}\bigr)s=\mathcal F^\Pi(Y,Z)\,s ,
 \label{eq:moving-curvature-def}
\end{equation}
where, as an operator on $\mathcal L_x$,
\begin{equation}
 \mathcal F^\Pi(Y,Z)=\Pi\,[Y\Pi,Z\Pi]\,\Pi=B_Y^*B_Z-B_Z^*B_Y .
 \label{eq:moving-curvature}
\end{equation}
The operator $\mathcal F^\Pi(Y,Z)$ is bounded and anti-Hermitian, and it depends only on $Y(x)$ and $Z(x)$. In coordinates,
\begin{equation}
 \mathcal F^\Pi_{\mu\nu}=[\nabla_\mu,\nabla_\nu]=B_\mu^*B_\nu-B_\nu^*B_\mu .
 \label{eq:moving-curvature-components}
\end{equation}
\end{proposition}
\begin{proof}
By the product rule $\nabla_Y\nabla_Zs=\Pi\,(Y\Pi)(Zs)+\Pi\,(YZs)$. Differentiating $s=\Pi s$ gives $Zs=(Z\Pi)s+\Pi(Zs)$, and the second term is annihilated by $\Pi(Y\Pi)$ because $\Pi(Y\Pi)\Pi=0$. Hence
\[
 \nabla_Y\nabla_Zs=\Pi\,(Y\Pi)(Z\Pi)\,s+\Pi\,(YZs).
\]
Subtracting the same expression with $Y$ and $Z$ exchanged, and using $YZs-ZYs=[Y,Z]s$, gives Eq.~\eqref{eq:moving-curvature-def} with $\mathcal F^\Pi(Y,Z)=\Pi[Y\Pi,Z\Pi]\Pi$. By Eq.~\eqref{eq:moving-block-forms} the upper left block of $(Y\Pi)(Z\Pi)$ is $B_Y^*B_Z$, which gives the second expression. Its adjoint is its negative. Coordinate fields commute, which gives Eq.~\eqref{eq:moving-curvature-components}.
\end{proof}

The product $B_Y^*B_Z=d_Y^*d_Z|_{\mathcal L_x}$ takes a vector of $\mathcal L_x$ into the complement along $Z$ and brings it back along $Y$. Call it the \emph{mixed return} of the two directions. For $Y=Z$ it is the positive return $B_Y^*B_Y$ of Section~\ref{sec:cohomology-index}. For two different directions it need not be self-adjoint, and Eq.~\eqref{eq:moving-curvature} says that the curvature is the difference of the two orders,
\begin{equation}
 \mathcal F^\Pi(Y,Z)=\bigl(d_Y^*d_Z-d_Z^*d_Y\bigr)\big|_{\mathcal L_x}.
 \label{eq:moving-mixed-return}
\end{equation}
Curvature is thus not the return of one map. It measures the failure of the two mixed returns to agree when the directions are exchanged, and it vanishes for a family with one parameter. For a family of rays, $r=1$, the mixed return is a number, the quantum geometric tensor of Provost and Vallee \cite{ProvostVallee1980}. Its real part is the natural metric on rays, pulled back to the parameters, and $\mathcal F^\Pi(Y,Z)=2i\operatorname{Im}(B_Y^*B_Z)$. For a subspace of finite dimension $r>1$ the mixed return is an operator. The scalar form
\begin{equation}
 \operatorname{Re}\Tr\bigl(B_Y^*B_Z\bigr)=\tfrac12\Tr\bigl((Y\Pi)(Z\Pi)\bigr)
 \label{eq:moving-scalar-form}
\end{equation}
is then symmetric and positive semidefinite in the two directions. It vanishes in a direction $Y$ exactly when $Y\Pi=0$, so it degenerates wherever a change of the parameters leaves the subspace fixed. For a subspace of infinite dimension the trace exists only under a further hypothesis, for instance that the tangent maps are Hilbert--Schmidt.

\subsection{Frames and transport around a loop}
\label{subsec:moving-frames}

Suppose that $\mathcal L_x$ has finite dimension $r$ on an open set $U\subseteq X$. A frame on $U$ is a map $\Xi:U\to\mathcal B(\C^r,\mathcal K)$ of class $C^2$ with
\begin{equation}
 \Xi^*\Xi=I_r,\qquad \Xi\,\Xi^*=\Pi .
 \label{eq:moving-frame}
\end{equation}
Frames exist locally. If the columns of $\Xi_0$ form an orthonormal basis of $\mathcal L_{x_0}$, then $\Xi(x)=\Pi(x)\,\Xi_0\,(\Xi_0^*\Pi(x)\Xi_0)^{-1/2}$ is a frame on the neighbourhood of $x_0$ where $\Xi_0^*\Pi(x)\Xi_0$ is invertible. In a frame every section is $s=\Xi c$ with $c:U\to\C^r$. Define the connection matrix and its curvature by
\begin{align}
 \mathcal A_\mu&=\Xi^*\partial_\mu\Xi ,
 \label{eq:moving-connection-matrix}\\
 \mathcal F^{(\Xi)}_{\mu\nu}&=\partial_\mu\mathcal A_\nu-\partial_\nu\mathcal A_\mu+[\mathcal A_\mu,\mathcal A_\nu].
 \label{eq:moving-frame-curvature}
\end{align}

\begin{proposition}[The connection in a frame]
\label{prop:moving-frame}
\textup{(a)} The connection matrix is anti-Hermitian, $\mathcal A_\mu^*=-\mathcal A_\mu$, and $\nabla_\mu(\Xi c)=\Xi\,(\partial_\mu c+\mathcal A_\mu c)$.

\textup{(b)} The matrix $\mathcal F^{(\Xi)}_{\mu\nu}$ represents the curvature of Proposition~\ref{prop:moving-curvature}:
\begin{equation}
 \mathcal F^{(\Xi)}_{\mu\nu}=\Xi^*\mathcal F^\Pi_{\mu\nu}\,\Xi=\Xi^*\bigl(B_\mu^*B_\nu-B_\nu^*B_\mu\bigr)\Xi .
 \label{eq:moving-frame-covariance}
\end{equation}

\textup{(c)} If $\Xi'=\Xi q$ with $q:U\to\mathrm U(r)$ of class $C^2$, then
\begin{equation}
 \mathcal A'_\mu=q^{-1}\mathcal A_\mu q+q^{-1}\partial_\mu q,\qquad
 \mathcal F^{(\Xi')}_{\mu\nu}=q^{-1}\mathcal F^{(\Xi)}_{\mu\nu}\,q .
 \label{eq:moving-change-of-frame}
\end{equation}
\end{proposition}
\begin{proof}
Differentiating $\Xi^*\Xi=I_r$ gives $(\partial_\mu\Xi)^*\Xi+\Xi^*\partial_\mu\Xi=0$, which is the first statement of (a). For the second, $\nabla_\mu(\Xi c)=\Xi\,\Xi^*\bigl((\partial_\mu\Xi)c+\Xi\,\partial_\mu c\bigr)=\Xi\,(\mathcal A_\mu c+\partial_\mu c)$. Thus $\nabla_\mu$ acts on coefficient vectors as $\partial_\mu+\mathcal A_\mu$, and for $c$ of class $C^2$ the commutator of two such operators is multiplication by the matrix in Eq.~\eqref{eq:moving-frame-curvature}. Comparing with Proposition~\ref{prop:moving-curvature} gives $\Xi\,\mathcal F^{(\Xi)}_{\mu\nu}c=\mathcal F^\Pi_{\mu\nu}\,\Xi c$, and multiplying by $\Xi^*$ gives (b). In (c), the connection matrix of $\Xi q$ is $q^*\Xi^*\bigl((\partial_\mu\Xi)q+\Xi\,\partial_\mu q\bigr)$ with $q^*=q^{-1}$, and the law for the curvature follows from (b) applied to both frames.
\end{proof}

Transport along a path does not need a frame. Let $\gamma:[0,1]\to X$ be a piecewise $C^1$ path and write $\dot\Pi=\frac{d}{dt}\Pi(\gamma(t))$. A section $s(t)\in\mathcal L_{\gamma(t)}$ along the path is called parallel when $\Pi\dot s=0$. Since $s=\Pi s$ gives $\dot s=\dot\Pi s+\Pi\dot s$, and $\Pi\dot\Pi\Pi=0$, a parallel section satisfies the linear equation
\begin{equation}
 \dot s=[\dot\Pi,\Pi]\,s .
 \label{eq:moving-transport}
\end{equation}
The generator $[\dot\Pi,\Pi]$ is bounded and anti-Hermitian, so this equation has a unitary propagator $T_\gamma(t)$ on $\mathcal K$. A direct calculation gives $[[\dot\Pi,\Pi],\Pi]=\dot\Pi$, hence $\frac{d}{dt}\bigl(T_\gamma^*\,\Pi(\gamma(t))\,T_\gamma\bigr)=0$ and
\begin{equation}
 T_\gamma(t)\,\Pi(\gamma(0))\,T_\gamma(t)^*=\Pi(\gamma(t)).
 \label{eq:moving-transport-intertwining}
\end{equation}
This is Kato's transformation function for a family of projections \cite{Kato}. Every solution of Eq.~\eqref{eq:moving-transport} that starts in $\mathcal L_{\gamma(0)}$ therefore stays in the moving subspace and is parallel. Transport carries $\mathcal L_{\gamma(0)}$ unitarily onto $\mathcal L_{\gamma(t)}$ for subspaces of any dimension. For a loop, the restriction of $T_\gamma(1)$ to $\mathcal L_{\gamma(0)}$ is a unitary of that fibre, the \emph{holonomy} $\operatorname{Hol}_\gamma$ of the loop. In a frame on a set containing the path, $s=\Xi c$ is parallel exactly when
\begin{equation}
 \dot c(t)+\mathcal A_\mu(\gamma(t))\,\dot\gamma^\mu(t)\,c(t)=0 ,
 \label{eq:moving-frame-transport}
\end{equation}
by part (a) of Proposition~\ref{prop:moving-frame}. The holonomy of a loop is then represented by a matrix in $\mathrm U(r)$, which a change of frame conjugates by $q(\gamma(0))$.

Holonomy in this sense should be kept apart from the sign in $\mathcal I^2=-I$ of Section~\ref{sec:neutral-graphs-phase}. That identity is algebraic. It holds for one fixed map on its active carrier and involves no path. Holonomy requires a connection and transport around an actual loop in the parameter space.

\subsection{Characteristic forms and integer flux}
\label{subsec:moving-characteristic}

Let $X$ now be a smooth manifold and $\Pi$ a smooth field of projections of finite rank $r$. Write $\mathcal L$ for the bundle of the subspaces $\mathcal L_x$. Collect the components of the curvature into the operator-valued two-form
\begin{equation}
 \mathcal F^\Pi=\tfrac12\,\mathcal F^\Pi_{\mu\nu}\,dx^\mu\wedge dx^\nu=\Pi\,(d\Pi\wedge d\Pi)\,\Pi .
 \label{eq:moving-curvature-form}
\end{equation}

\begin{proposition}[Characteristic forms of a moving subspace]
\label{prop:moving-chern}
For every $k\ge1$ the form
\begin{equation}
 \operatorname{ch}_k(\mathcal L,\nabla)
 =\frac1{k!}\Bigl(\frac{i}{2\pi}\Bigr)^k\Tr_{\mathcal L}\bigl((\mathcal F^\Pi)^{\wedge k}\bigr)
 =\frac1{k!}\Bigl(\frac{i}{2\pi}\Bigr)^k\Tr\bigl(\Pi\,(d\Pi)^{\wedge 2k}\bigr)
 \label{eq:moving-chern-character}
\end{equation}
is closed. Its de Rham class does not depend on the connection chosen on $\mathcal L$, and it is the component of degree $2k$ of the Chern character of $\mathcal L$, whose component of degree zero is the rank \cite{KobayashiNomizu,Husemoller}. In particular the first Chern form and its components are
\begin{equation}
 c_1(\mathcal L,\nabla)=\frac{i}{2\pi}\Tr_{\mathcal L}\mathcal F^\Pi,\qquad
 c_1(\mathcal L,\nabla)_{\mu\nu}=\frac{i}{2\pi}\Tr_{\mathcal L}\bigl(B_\mu^*B_\nu-B_\nu^*B_\mu\bigr),
 \label{eq:moving-first-chern}
\end{equation}
and
\begin{equation}
 \operatorname{ch}_2(\mathcal L,\nabla)=\frac12\Bigl(\frac{i}{2\pi}\Bigr)^2\Tr_{\mathcal L}\bigl(\mathcal F^\Pi\wedge\mathcal F^\Pi\bigr).
 \label{eq:moving-second-chern}
\end{equation}
\end{proposition}
\begin{proof}
By Proposition~\ref{prop:moving-tangent} each $\partial_\mu\Pi$ anticommutes with $J_\Pi$. A product of two of them commutes with $J_\Pi$ and so with $\Pi$. Hence $(\mathcal F^\Pi)^{\wedge k}=\Pi\,(d\Pi)^{\wedge 2k}$, which gives the second expression. All operators here have finite rank. Since $d(d\Pi)=0$, the exterior derivative of $\Tr(\Pi\,(d\Pi)^{\wedge 2k})$ is $\Tr((d\Pi)^{\wedge(2k+1)})$. Each coefficient of that form is the trace of a product $T$ of $2k+1$ operators that anticommute with $J_\Pi$. Such a product satisfies $J_\Pi TJ_\Pi=-T$, so $\Tr T=\Tr(J_\Pi TJ_\Pi)=-\Tr T=0$. This proves closedness. The statements about the cohomology class are the Chern--Weil theorem for the unitary connection $\nabla$ on the Hermitian bundle $\mathcal L$. The displayed cases are $k=1$ and $k=2$ with Eq.~\eqref{eq:moving-curvature-components}.
\end{proof}

\begin{corollary}[Integer flux]
\label{cor:moving-flux}
Let $\mathsf Z$ be a closed oriented surface and $f:\mathsf Z\to X$ a smooth map. Then
\begin{equation}
 \frac{i}{2\pi}\int_{\mathsf Z}f^*\,\Tr_{\mathcal L}\mathcal F^\Pi\in\mathbb Z .
 \label{eq:moving-flux}
\end{equation}
The integer is the pairing of the first Chern class of $\mathcal L$ with the class of the cycle. For a line bundle, $r=1$, the trace may be omitted.
\end{corollary}
\begin{proof}
The field $\Pi\circ f$ on $\mathsf Z$ satisfies the hypotheses, so assume $X=\mathsf Z$. Choose a triangulation of $\mathsf Z$ fine enough that every closed triangle $\Delta$ lies in the domain of a frame $\Xi_\Delta$, with connection matrix $\mathcal A^{(\Delta)}$. On that domain $\Tr_{\mathcal L}\mathcal F^\Pi=\Tr\mathcal F^{(\Xi_\Delta)}=d\,\Tr\mathcal A^{(\Delta)}$ by Proposition~\ref{prop:moving-frame}, since the trace of a commutator vanishes. Stokes' theorem gives $\int_\Delta\Tr_{\mathcal L}\mathcal F^\Pi=\oint_{\partial\Delta}\Tr\mathcal A^{(\Delta)}$.

Along an oriented edge $e$ of $\Delta$ from a vertex $v$ to a vertex $v'$, let $P_e$ be the solution matrix of Eq.~\eqref{eq:moving-frame-transport}. It is unitary with $\det P_e=\exp(-\int_e\Tr\mathcal A^{(\Delta)})$, and the transport of Eq.~\eqref{eq:moving-transport} from $\mathcal L_v$ to $\mathcal L_{v'}$ is $\tau_e=\Xi_\Delta(v')P_e\,\Xi_\Delta(v)^*$. Fix for each vertex an isometry $\beta_v:\C^r\to\mathcal K$ with range $\mathcal L_v$, and put $\delta_e=\det(\beta_{v'}^*\tau_e\beta_v)$. This number has modulus one, does not refer to $\Delta$, and is inverted when $e$ is reversed. With $\kappa_\Delta(v)=\det(\Xi_\Delta(v)^*\beta_v)$ one finds $\delta_e=\kappa_\Delta(v')^{-1}\det P_e\,\kappa_\Delta(v)$. Around the boundary of $\Delta$ the factors $\kappa_\Delta$ cancel, so
\[
 \prod_{e\subset\partial\Delta}\delta_e=\exp\Bigl(-\oint_{\partial\Delta}\Tr\mathcal A^{(\Delta)}\Bigr)
   =\exp\Bigl(-\int_\Delta\Tr_{\mathcal L}\mathcal F^\Pi\Bigr).
\]
In the product over all triangles each edge occurs twice, with opposite orientations, and the left sides multiply to one. Therefore $\exp(-\int_{\mathsf Z}\Tr_{\mathcal L}\mathcal F^\Pi)=1$, which is Eq.~\eqref{eq:moving-flux}. The identification with the Chern class is part of the Chern--Weil theorem.
\end{proof}

The argument applies equally to a smooth singular two-cycle with integer coefficients. After subdivision each simplex lies in the domain of a frame, and the edge factors cancel because the boundary of the cycle vanishes.

The convention here is anti-Hermitian. In terms of the Hermitian two-form $\mathcal F_{\mathrm H}=i\mathcal F^\Pi$ the corollary reads $(2\pi)^{-1}\int\Tr\mathcal F_{\mathrm H}\in\mathbb Z$.

The integer belongs to a closed surface. The holonomy of a loop is not quantized. For a loop $\gamma$ that bounds a disc $\mathsf S$ inside the domain of a frame, Stokes' theorem and Eq.~\eqref{eq:moving-frame-transport} give
\begin{equation}
 \det\operatorname{Hol}_\gamma=\exp\Bigl(-\int_{\mathsf S}\Tr_{\mathcal L}\mathcal F^\Pi\Bigr),
 \label{eq:moving-holonomy-flux}
\end{equation}
which for $r=1$ is the holonomy itself. This can be any number of modulus one, also for a contractible loop. Curvature and holonomy are different things. Curvature is defined at a point and needs two directions. Holonomy belongs to a loop, and Eq.~\eqref{eq:moving-holonomy-flux} relates it to the curvature on a surface that the loop bounds, not to the family restricted to the loop.

Because the curvature is the antisymmetrized mixed return, every characteristic form is a trace polynomial in antisymmetrized mixed returns, as Eq.~\eqref{eq:moving-first-chern} shows in the first case. These are exact statements about the bundle of subspaces. Identifying a characteristic number with a physical charge requires additional data of a realization and is not part of them.

\subsection{Turning and waiting in the pair model}
\label{subsec:moving-pair}

The pair model of Section~\ref{sec:preparation-effective-dynamics} supplies a family with two parameters. The moving subspace is the prepared plane inside the six-dimensional output space. The splitting into $\Psi$ and $\Phi$ does not move. Let $P_\theta=A_\theta A_\theta^*$ project onto the prepared plane $\Span\{g,b\}$, with $b=u_\theta e+v_\theta m$ and $d=-v_\theta e+u_\theta m$, and let $U(t)=\exp(-itH_s/\hbar)$ with $\delta=j-s>0$. Preparing at the setting $\theta$ and waiting for a time $t$ gives the plane with projection
\begin{equation}
 \Pi(\theta,t)=U(t)\,P_\theta\,U(t)^* .
 \label{eq:moving-pair-family}
\end{equation}
At $t=0$ the two tangent maps are
\begin{equation}
 B_\theta=2\ket{d}\bra{b},\qquad
 B_t=-\frac{i}{\hbar}\,(I-P_\theta)H_sP_\theta=\frac{i\delta}{\hbar}\,u_\theta v_\theta\ket{d}\bra{b},
 \label{eq:moving-pair-tangents}
\end{equation}
since $\partial_\theta b=2d$ and $\partial_t\Pi=-(i/\hbar)[H_s,\Pi]$. For any bounded Hamiltonian and any projection, the tangent map of waiting is $-i/\hbar$ times the coupling $C$ of Eq.~\eqref{eq:prepared-schur}. Turning and waiting both take $b$ into the complex line of $d$, and $B_t=(i\delta u_\theta v_\theta/2\hbar)\,B_\theta$. The factor is imaginary. For $0<\theta<\pi/4$ the two are therefore orthogonal as real tangent directions, $\operatorname{Re}\Tr(B_\theta^*B_t)=0$ in Eq.~\eqref{eq:moving-scalar-form}, and their mixed return $B_\theta^*B_t=(2i\delta/\hbar)u_\theta v_\theta\ket{b}\bra{b}$ is anti-Hermitian. That is how two maps into the same complex line produce curvature. By Eq.~\eqref{eq:moving-curvature-components},
\begin{equation}
 \mathcal F^\Pi_{\theta t}=B_\theta^*B_t-B_t^*B_\theta=\frac{4i\delta}{\hbar}\,u_\theta v_\theta\ket{b}\bra{b}
 \qquad\text{at }t=0 ,
 \label{eq:moving-pair-curvature}
\end{equation}
and at other times the same operator conjugated by $U(t)$. Its trace is $(2i\delta/\hbar)\sin(4\theta)$.

Which closed surface carries an integer? The settings $0\le\theta\le\pi/4$ and the waiting times modulo the period $t_\delta=2\pi\hbar/\delta$ form a cylinder. At $\theta=0$ and at $\theta=\pi/4$ the prepared plane is $\Span\{g,e\}$ and $\Span\{g,m\}$ for every waiting time, because $e$ and $m$ are eigenvectors. Collapsing each of the two end circles to a point gives a sphere, on which every line of $\Span\{e,m\}$ occurs once as the line of $U(t)b$ and $\Pi$ is a smooth field of distinct planes. With the orientation $d\theta\wedge dt$,
\begin{equation}
 \frac{i}{2\pi}\int_0^{\pi/4}\!\!\int_0^{t_\delta}\Tr\mathcal F^\Pi_{\theta t}\,dt\,d\theta
   =-4\int_0^{\pi/4}u_\theta v_\theta\,d\theta=-1 ,
 \label{eq:moving-pair-flux}
\end{equation}
the integer of Corollary~\ref{cor:moving-flux} for this sphere. The plane itself has period $\pi/2$ in $\theta$. Over $0\le\theta\le\pi/2$ the parameters form a torus, the second half covers the same sphere with the opposite orientation, and the integral is zero.

Now fix $\theta$ and let $t$ run over one period. This loop is a family with one parameter and has no curvature of its own. For $0<\theta<\pi/4$ its holonomy is nevertheless nontrivial. Transport around the loop fixes $g$ and multiplies $b$ by
\begin{equation}
 e^{2\pi iu_\theta^2} ,
 \label{eq:moving-pair-holonomy}
\end{equation}
because the frame vector $\psi(t)=U(t)b$ has $\langle\psi,\dot\psi\rangle=-i\delta u_\theta^2/\hbar$. By Eq.~\eqref{eq:moving-holonomy-flux} this is $\exp(-2\pi iv_\theta^2)$, where $2\pi iv_\theta^2$ is the integral of the trace of the curvature over the cap of the sphere between the setting $0$ and the setting $\theta$. The excitation itself returns exactly, $U(t_\delta)b=b$. The two statements differ by the factor $\exp(-(i/\hbar)\int_0^{t_\delta}\langle\psi,H_s\psi\rangle\,dt)=e^{-2\pi iu_\theta^2}$ contributed by the mean energy $\delta u_\theta^2$. The phase in Eq.~\eqref{eq:moving-pair-holonomy} is thus the geometric phase of the cyclic evolution of the excitation \cite{AharonovAnandan1987}. It is fixed by the weight $u_\theta^2$ that the response of Section~\ref{sec:preparation-effective-dynamics} assigns to the line at energy $\delta$. At $\theta=\pi/8$ the two line weights are equal, and transport around one period acts by $+1$ on the ground and by $-1$ on the excitation.
